% ECONOMICS_PROBLEM: {"id": "L51-508", "title": "Reputation versus public regulation", "jel_code": "L51", "source_id": "OP-0405"}
% FIRST_POSTED: 2026-10-09
% ATTEMPT_STATUS: partial
% ENTRY_KIND: research_attempt
\documentclass[11pt,a4paper]{article}
\usepackage[T1]{fontenc}
\usepackage[utf8]{inputenc}
\usepackage[margin=27mm]{geometry}
\usepackage{amsmath,amssymb,amsthm,mathtools}
\usepackage[hidelinks]{hyperref}
\setlength{\parindent}{0pt}
\setlength{\parskip}{5pt}
\newtheorem{theorem}{Theorem}[section]
\newtheorem{proposition}[theorem]{Proposition}
\newtheorem{lemma}[theorem]{Lemma}
\newcommand{\R}{\mathbb R}
\newcommand{\E}{\mathbb E}
\DeclareMathOperator{\Var}{Var}
\DeclareMathOperator{\Cov}{Cov}
\begin{document}
\section*{Reputation and regulation when safety is incompletely visible}
\section{Partial result and the target population}
The comparison concerns all households needing a repair, including those who cannot hire a provider. Ratings among completed jobs do not define that population. This attempt gives a hidden-quality incentive model, derives the least inspection intensity needed for a safety constraint in that model, and supplies sharp audit-selection bounds. It leaves the actual three-year regime ranking and endogenous provider composition unresolved.

Let a provider choose safety effort $e\in[0,1]$, with real cost $ke^2/2$, $k>0$. The probability of a harmful defect is $1-e$, and the harm costs the household $\ell>0$ in the declared money units. Suppose the expected discounted private reward generated by safety-sensitive reputation is $ve$, where $v\ge0$. An independent inspection detects a defect with probability $a\in[0,1]$ and imposes a collected fine $F\ge0$ when a defect is detected. The provider's expected payoff, up to terms independent of effort, is
\[
 ve-\tfrac12ke^2-aF(1-e).
\]
The reward $v$ is an explicit reduced-form continuation-value primitive. A review system focused entirely on politeness or visible finish can have $v=0$ even if consumers strongly value safety.

\section{Exact effort and a safety feasibility threshold}
Strict concavity gives the unique effort
\[
 e(a)=\operatorname{clip}_{[0,1]}((v+aF)/k). \tag{1}
\]
At an interior optimum the derivative is $v+aF-ke=0$; checking the endpoint signs gives the clipping rule. Suppose the regime must achieve defect probability at most $\rho\in[0,1]$. For $F>0$, define
\[
 a_{\min}=\max\{0,[k(1-\rho)-v]/F\}. \tag{2}
\]
The safety target is feasible exactly when $a_{\min}\le1$, and then it is met exactly for $a\ge a_{\min}$. Indeed the target is equivalent to $e(a)\ge1-\rho$, which is equivalent to $v+aF\ge k(1-\rho)$ for the stated ranges. With $F=0$, feasibility instead requires $v\ge k(1-\rho)$. These statements include the endpoints $\rho=0$ and $\rho=1$.

Reputation alone suffices in this benchmark precisely when $v\ge k(1-\rho)$. Perfectly informative reviews on an unrelated quality dimension do not establish that inequality. When $v=0$, reputational incentives induce zero safety effort unless monitoring supplies a private incentive. Conversely, sufficiently large $v$ can make inspections unnecessary for the given constraint. These are conditional possibilities, not a universal substitution or complementarity claim.

For example $k=4,v=1,F=5,\rho=0.25$ gives $a_{\min}=0.4$, which induces $e=0.75$. If reputation improves to $v=2$, required inspection falls to $0.2$. But a shift toward harder-to-monitor providers can increase $k$ or lower the effective collected fine, reversing this comparison. The composition of providers is held fixed in (1), and cannot be assumed fixed in the catalogue's market-level evaluation.

\section{Real welfare, fines, and an interior monitoring optimum}
Assume the household service benefit before defects is fixed, and inspection has real cost $ca$ per eligible household, $c\ge0$. Under equal social weights with fines rebated inside the evaluated economy, fines are transfers and the real loss is
\[
 L(a)=\tfrac12k e(a)^2+\ell[1-e(a)]+ca. \tag{3}
\]
The reputational reward is also taken to be a transfer or continuation-rent reallocation within this restricted calculation; if producing it consumes resources, those costs must be added. On the interior region $0<v+aF<k$,
\[
 L'(a)=F(e(a)-\ell/k)+c,\qquad L''(a)=F^2/k>0. \tag{4}
\]
Thus an interior unconstrained optimum has $e=\ell/k-c/F$, when this expression is feasible and $F>0$. For a safety-constrained optimum, examine that stationary point, the effort-clipping breakpoints, and the feasible endpoints $a_{\min},1$; on each piece the function is convex quadratic or affine. This finite comparison is exact. It is incorrect simply to add collected fines to social costs while also subtracting them from provider profits, since their recipients then disappear from the account.

Consumer net welfare differs from (3): provider profits and transfers have different incidence. For two regimes, a lower repair payment can benefit consumers and reduce providers' income without itself saving real resources. The catalogue's distinction between $W_C$ and $W_T$ is therefore essential. A comparison using one objective cannot establish a ranking for the other without the missing profit and resource-cost terms.

\section{Selective inspections and sharp harm bounds}
Let $A$ indicate whether a household receives a completed independent defect audit. Suppose the audited harm is $H\in[0,H_{\max}]$, audit probability is $p$, and observed mean harm is $m$. Without audit-selection restrictions,
\[
 E[H]\in[pm,\quad pm+(1-p)H_{\max}]. \tag{5}
\]
For the tail event $T=\mathbf1\{H>h_0\}$ with observed audited rate $t$,
\[
 \Pr(H>h_0)\in[pt,\quad pt+1-p]. \tag{6}
\]
Assigning missing harms zero or $H_{\max}$ establishes (5); assigning missing tail indicators zero or one establishes (6), assuming $0\le h_0<H_{\max}$. Both endpoints can be attained with compatible outcome laws. Joint mean-tail bounds need their shared distribution and are not generally the rectangle formed by (5)--(6). For $h_0$ outside the harm support the tail event is constant, and (6) should be replaced by that exact value.

A safety cap $\Pr(H>h_0)\le\rho$ is certified from these data without selection assumptions only if $pt+1-p\le\rho$. Even zero observed defects among a small audited fraction can fail that criterion. Voluntary reviews have an additional reporting-selection mechanism and may not measure latent defects at all.

If audits are genuinely randomized independently of harm among the full eligible population, complete audit follow-up identifies the harm distribution from audited households. With $n$ independent audited households, an upper Hoeffding bound for the tail rate is
$\min\{1,\widehat t+\sqrt{\log(1/\alpha)/(2n)}\}$. This is a finite-sample certificate under independent sampling, not under repeated inspections sharing an unmodeled provider shock. Provider-level clustering requires an appropriate independent-unit design. Zero audits give no nontrivial certificate.

\section{Endogenous entry and no-service outcomes}
A licensing rule can remove unsafe providers and also deter safe low-margin providers. In a minimal accounting example, take eligible-household service probability $s_r$, service value $b_r$ net of real production cost and expected harm, and no-service value zero. Ignoring other costs, total benefit is $s_rb_r$. A regime improving conditional service value from ten to twelve but lowering access from one to one-half reduces total benefit from ten to six. The example demonstrates the denominator issue without assuming that access loss is always dominant.

For observed data, count unfilled requests and offline substitutions in the eligible population, and specify their outcomes. Endogenous prices, reputation histories and provider effort can then enter a market model or a market-level assignment comparison. The incentive threshold (2) alone cannot establish which policy changes provider composition favorably.

\section{Remaining gap}
The benchmark locates the information dimension on which reputation must operate, and the audit bounds show what unobserved defects can conceal. Actual regime rankings need representative request-level data, independent latent-defect audits, provider profits and real monitoring costs. Long-run entry and selective platform participation are not solved by the one-provider effort model. The full catalogue question therefore remains partial.

\section{Completion Estimate}
50\%

This is a post hoc, heuristic estimate for the full catalogue target.
The attempt derives safety-effort and minimum-inspection thresholds, correct transfer accounting, and sharp selective-audit harm bounds with safety certificates. Full reputation-regulation rankings still need representative repair demand, latent-defect audits, provider entry and platform participation responses, and consumer and producer outcome data over three years.

\begin{thebibliography}{9}
\bibitem{f} C. Farronato (2025), \emph{Digital Platforms as Information Aggregators: Implications for their Design and Regulation}, conference draft, Section 2.1. \url{https://conference.nber.org/confer/2025/DTs25/farronato.pdf}. Source motivation for combined platform monitoring and public regulation; the incentive and audit calculations above are conditional derivations.
\end{thebibliography}

\end{document}
